\makeatletter
\def\input@path{{preprints-template/preprints_template_latex/}}
\let\preprint@bibliographystyle\bibliographystyle
\renewcommand{\bibliographystyle}[1]{%
  \preprint@bibliographystyle{preprints-template/preprints_template_latex/#1}}
\makeatother

\documentclass[preprints,article,submit,oneauthor]{Definitions/mdpi}
\setlength{\headheight}{18pt}

\makeatletter
\let\bibliographystyle\preprint@bibliographystyle
\makeatother

\graphicspath{{preprints-template/preprints_template_latex/}}

\firstpage{1}
\makeatletter
\setcounter{page}{\@firstpage}
\makeatother
\pubvolume{1}
\issuenum{1}
\articlenumber{0}
\pubyear{2026}
\copyrightyear{2026}
\datereceived{}
\daterevised{}
\dateaccepted{}
\datepublished{}

\newcounter{manuscriptresult}[section]
\renewcommand{\themanuscriptresult}{\thesection.\arabic{manuscriptresult}}
\theoremstyle{mdpi}
\newtheorem{theorem}[manuscriptresult]{Theorem}
\newtheorem{lemma}[manuscriptresult]{Lemma}
\newtheorem{proposition}[manuscriptresult]{Proposition}
\newtheorem{corollary}[manuscriptresult]{Corollary}

\theoremstyle{definition}
\newtheorem{remark}[manuscriptresult]{Remark}

\newcommand{\C}{\mathbb C}
\newcommand{\D}{\mathbb D}
\newcommand{\Mn}{\mathbb M_n(\C)}

\Title{The sharp universal spectral constant for scaled
\texorpdfstring{$q$}{q}-numerical ranges}
\Author{Shanmu Jin}
\AuthorNames{Shanmu Jin}
\address{Department of Neurosurgery, Peking Union Medical College
  Hospital, No.~1 Shuaifuyuan, Dongcheng District, Beijing 100730,
  China\\
jinshanmu1129@pku.edu.cn}
\corres{Correspondence: jinshanmu1129@pku.edu.cn}

\abstract{\texorpdfstring{%
Let $A$ be an $n\times n$ complex matrix, $n\geq2$, and let
$q\in\C$ satisfy $0<\lvert q\rvert\leq1$.  The $q$-numerical range
and its scaled form are
\[
 W_q(A)=\{y^*Ax:\lVert x\rVert_2=\lVert y\rVert_2=1,\ y^*x=q\},
 \qquad
 \Omega_q(A)=q^{-1}W_q(A).
\]
We prove that the sharp universal spectral constant for the family of
scaled $q$-numerical ranges is
\[
 C_q=\max\left\{1,
 \frac{2\lvert q\rvert}
 {1+\sqrt{1-\lvert q\rvert^2}}\right\}.
\]
Thus, for every rational function $f$ pole-free on $\Omega_q(A)$,
$\lVert f(A)\rVert$ is at most $C_q$ times the maximum modulus of $f$
on $\Omega_q(A)$, and $C_q$ is the least constant valid uniformly in
$n$ and $A$.
The proof is a transfer from the sharp numerical-range theorem.  Its
geometric input shows that the numerical range of every adapted
rank-one stretch of $A$ lies in the scaled $q$-range.  Applying a
numerical-range spectral-set estimate to these stretches produces a
simultaneous family of Schur estimates.  A disk automorphism adapted to
a maximizing singular-vector pair then extracts a norm bound for the
original matrix.  More generally, a universal numerical-range
spectral constant $K$ transfers to the constant
$\max\{1,K/\kappa(q)\}$ on the scaled $q$-range, where
$\kappa(q)=(1+\sqrt{1-\lvert q\rvert^2})/\lvert q\rvert$.  The two
lower-bound branches defining $C_q$ are attained, respectively, by the
constant rational function and by the identity polynomial on a
$2\times2$ square-zero matrix.
}{%
For a complex matrix A and a parameter q satisfying 0<|q|<=1, we
determine the sharp universal spectral constant for the scaled
q-numerical range in the rational spectral-set formulation.  The proof
applies the numerical-range spectral-set estimate to adapted rank-one
stretches and then uses a disk automorphism chosen from a maximizing
singular-vector pair.  More generally, every universal numerical-range
spectral constant transfers explicitly to scaled q-numerical ranges.
Constant functions and a two-dimensional square-zero matrix attain the
two lower-bound branches, respectively.
}}

\keyword{q-numerical range; spectral set; Crouzeix conjecture;
numerical range; rank-one similarity; Schur function}
\MSC{Primary 47A25; Secondary 47A12, 15A60}

\begin{document}

\section{Introduction}
\label{sec:introduction}

For $A\in\Mn$, its numerical range is
\[
 W(A)=\{x^*Ax:x\in\C^n,\ \lVert x\rVert_2=1\}.
\]
It is compact and convex, and it contains the spectrum of $A$.
For vectors, $\lVert\cdot\rVert_2$ denotes the Euclidean norm.  For a
matrix $T$, the unsubscripted norm is the induced Euclidean operator norm,
\[
 \lVert T\rVert:=
 \max_{\lVert x\rVert_2=1}\lVert Tx\rVert_2.
\]
For $C\geq1$, a compact set $E\subset\C$ containing $\sigma(A)$ is a
$C$-spectral set for $A$ if
\begin{equation}
 \lVert g(A)\rVert
 \leq C\max_{z\in E}\lvert g(z)\rvert
 \label{eq:spectral-set-definition}
\end{equation}
for every rational function $g$ pole-free on $E$.
Crouzeix conjectured that $W(A)$ is always a $2$-spectral set
\cite{Crouzeix2004}, and the sharp constant $2$ was established in
\cite{Jin2026}.

Throughout, $q\in\C$.  For $n\geq2$ and $0<\lvert q\rvert\leq1$, the
$q$-numerical range is
\begin{equation}
 W_q(A)=
 \{y^*Ax:\lVert x\rVert_2=\lVert y\rVert_2=1,\ y^*x=q\},
 \qquad
 \Omega_q(A)=\frac1qW_q(A).
 \label{eq:q-range-definition}
\end{equation}
The set $W_q(A)$ is compact, and Tsing proved its convexity
\cite{Tsing1984}.  The scaled set $\Omega_q(A)$ contains $W(A)$ and
hence $\sigma(A)$.  O'Loughlin and Rani formulated the polynomial
inequality below as Conjecture 4.2 \cite{OLoughlinRani2026}.  The
theorem proves that conjecture together with its rational spectral-set
extension.

\begin{theorem}[Main theorem]
\label{thm:main}
Let $n\geq2$, let $A\in\Mn$, let $0<\lvert q\rvert\leq1$, and let
$f$ be a rational function pole-free on $\Omega_q(A)$.  Then
\begin{equation}
 \lVert f(A)\rVert
 \leq
 C_q\max_{z\in\Omega_q(A)}\lvert f(z)\rvert,
 \qquad
 C_q=
 \max\left\{1,
 \frac{2\lvert q\rvert}
 {1+\sqrt{1-\lvert q\rvert^2}}\right\}.
 \label{eq:main-bound}
\end{equation}
Equivalently, $\Omega_q(A)$ is a $C_q$-spectral set for $A$.
In particular, the same inequality holds for every polynomial $f$.
\end{theorem}

The proof has three steps.  First, Section~\ref{sec:similarity-geometry}
proves nesting of the scaled ranges and shows that every adapted
rank-one stretch has numerical range inside $\Omega_q(A)$.  Second,
Section~\ref{sec:orbit-extraction} proves that uniform Schur estimates
over precisely these stretches force an improved norm estimate for the
original matrix.  Third, Section~\ref{sec:main-proof} packages the two
facts into a transfer principle for any universal numerical-range
spectral constant $K$.  Taking $K=2$ gives \eqref{eq:main-bound}; a
$2\times2$ nilpotent matrix proves sharpness.

\section{Rank-one stretch geometry}
\label{sec:similarity-geometry}

We use the inner product
\[
 \langle x,y\rangle=y^*x,
\]
which is linear in its first variable.  Thus
$W_q(A)=\{\langle Ax,y\rangle:\lVert x\rVert_2=\lVert y\rVert_2=1,\
\langle x,y\rangle=q\}$.
Changing the phase of $y$ in \eqref{eq:q-range-definition} gives
\begin{equation}
 \Omega_q(A)=\Omega_{\lvert q\rvert}(A).
 \label{eq:phase-reduction}
\end{equation}
We may therefore put
\begin{equation}
 r=\lvert q\rvert,\qquad
 \tau=\sqrt{1-r^2},\qquad
 \kappa=\frac{1+\tau}{r}.
 \label{eq:parameter-normalization}
\end{equation}
The equivalent identities
\begin{equation}
 r=\frac{2\kappa}{\kappa^2+1},
 \qquad
 \frac\tau r=\frac{\kappa^2-1}{2\kappa},
 \qquad
 \frac2\kappa=\frac{2r}{1+\sqrt{1-r^2}}
 \label{eq:kappa-identities}
\end{equation}
will be used repeatedly.

For $0<r<1$, division by $r$ in Tsing's Lemma 5 gives the disk formula
below \cite{Tsing1984}.  The endpoint $r=1$ is included in the proof.
The formula will also be used in the sharpness argument.  For a unit
vector $x$, put
\[
 \rho_A(x)=
 \left(\lVert Ax\rVert_2^2-
       \lvert\langle Ax,x\rangle\rvert^2\right)^{1/2}.
\]
The radicand is nonnegative by Cauchy--Schwarz.

\begin{lemma}[Disk formula and nesting]
\label{lem:q-range-nesting}
Let $A\in\Mn$, $n\geq2$, and $0<r\leq1$.  Then
\begin{equation}
 \Omega_r(A)=
 \bigcup_{\lVert x\rVert_2=1}
 \left\{z\in\C:
  \lvert z-\langle Ax,x\rangle\rvert
  \leq\frac{\sqrt{1-r^2}}r\rho_A(x)\right\}.
 \label{eq:q-disk-formula}
\end{equation}
Consequently, if $0<r\leq s\leq1$, then
\[
 \Omega_s(A)\subseteq\Omega_r(A).
\]
In particular, $W(A)=\Omega_1(A)\subseteq\Omega_r(A)$.
\end{lemma}

\begin{proof}
Let $x,y$ be unit vectors with $\langle x,y\rangle=r$, and put
$e_x=Ax-\langle Ax,x\rangle x$.  Then $e_x\perp x$ and
$\lVert e_x\rVert_2=\rho_A(x)$.  Since
$y-rx\perp x$ and $\lVert y-rx\rVert_2=\sqrt{1-r^2}$,
Cauchy--Schwarz gives
\[
 \left\lvert
  \frac1r\langle Ax,y\rangle-\langle Ax,x\rangle
 \right\rvert
 =\frac1r\lvert\langle e_x,y-rx\rangle\rvert
 \leq\frac{\sqrt{1-r^2}}r\rho_A(x).
\]
This proves the inclusion from left to right in
\eqref{eq:q-disk-formula}.

Conversely, suppose that $\zeta$ belongs to the disk indexed by a unit
vector $x_0$.  On the unit sphere define
\[
 G(x)=\lvert\zeta-\langle Ax,x\rangle\rvert
      -\frac{\sqrt{1-r^2}}r\rho_A(x).
\]
The function $G$ is continuous and $G(x_0)\leq0$.  A complex matrix has
a unit eigenvector $e$, and $\rho_A(e)=0$, so $G(e)\geq0$.  The unit
sphere in $\C^n$ is path connected for $n\geq2$; hence there is a unit
vector $u$ with $G(u)=0$.

Write $\alpha=\langle Au,u\rangle$ and $d=\rho_A(u)$.  Choose a unit
vector $w\perp u$ such that $\langle Au,w\rangle=d$: when $d>0$, take
$w=(Au-\alpha u)/d$, and when $d=0$, take any unit vector orthogonal to
$u$.  If $d>0$, set
\[
 \beta=\frac{r\,\overline{\zeta-\alpha}}d;
\]
if $d=0$, set $\beta=\sqrt{1-r^2}$.  The equality $G(u)=0$ gives
$\lvert\beta\rvert=\sqrt{1-r^2}$.  Thus
$y=ru+\beta w$ is a unit vector with $\langle u,y\rangle=r$, and
if $d>0$, the definition of $\beta$ gives
$\overline\beta d/r=\zeta-\alpha$.  If $d=0$, then $G(u)=0$ gives
$\zeta=\alpha$.  In either case,
\[
 \frac1r\langle Au,y\rangle
 =\alpha+\frac{\overline\beta d}{r}=\zeta.
\]
This proves the reverse inclusion.

Finally, the coefficient $\sqrt{1-r^2}/r$ decreases with $r$ on
$(0,1]$.  Comparing the disks in \eqref{eq:q-disk-formula} proves
nesting, and the case $r=1$ gives $\Omega_1(A)=W(A)$.
\end{proof}

For a unit vector $v$, let $P_v$ be the orthogonal projection onto
$\C v$ and define the rank-one stretch
\begin{equation}
 S_v=I+(\kappa-1)P_v.
 \label{eq:rank-one-stretch}
\end{equation}

\begin{lemma}[Rank-one stretch geometry]
\label{lem:rank-one-stretch}
Let $A\in\Mn$, $n\geq2$, and $0<r\leq1$, with $\kappa$ as in
\eqref{eq:parameter-normalization}.  For every unit vector $v$,
$S_v\succ0$,
\[
 S_v^{-1}=I+(\kappa^{-1}-1)P_v,
\]
and
\begin{equation}
 W(S_vAS_v^{-1})\subseteq\Omega_r(A).
 \label{eq:rank-one-stretch-containment}
\end{equation}
\end{lemma}

\begin{proof}
The formulas for $S_v$ and $S_v^{-1}$ follow from
$P_v^2=P_v$; positivity follows from $\kappa\geq1$.  Fix a unit vector
$u$ and put
\[
 \ell=\lVert S_v^{-1}u\rVert_2,\qquad
 m=\lVert S_vu\rVert_2,
 \qquad \theta=\lVert P_vu\rVert_2^2.
\]
Orthogonality of $P_vu$ and $(I-P_v)u$ gives
\[
 m^2=\kappa^2\theta+1-\theta,
 \qquad
 \ell^2=\kappa^{-2}\theta+1-\theta.
\]
Since $0\leq\theta\leq1$,
\[
 \ell^2m^2
 =1+(\kappa-\kappa^{-1})^2\theta(1-\theta)
 \leq\frac14(\kappa+\kappa^{-1})^2.
\]
Moreover,
$1=\lvert\langle S_v^{-1}u,S_vu\rangle\rvert\leq \ell m$.
Because $S_v$ is invertible, $\ell,m>0$.  Therefore, with
$s=(\ell m)^{-1}$,
\[
 r\leq s\leq1,
 \qquad
 \frac12(\kappa+\kappa^{-1})=\frac1r.
\]
The vectors
\[
 x=\frac{S_v^{-1}u}{\ell},
 \qquad
 y=\frac{S_vu}{m}
\]
are unit vectors with $\langle x,y\rangle=s$, and
\[
 \langle S_vAS_v^{-1}u,u\rangle
 =\frac1s\langle Ax,y\rangle
 \in\Omega_s(A)\subseteq\Omega_r(A)
\]
by Lemma~\ref{lem:q-range-nesting}.  This proves
\eqref{eq:rank-one-stretch-containment}.
\end{proof}

\section{Rank-one stretch extraction}
\label{sec:orbit-extraction}

Write $\D=\{z\in\C:\lvert z\rvert<1\}$ for the open unit disk.
A convenient family of disk automorphisms is
\[
 \varphi_{\omega,a}(z)=\frac{\omega z-a}{1-a\omega z},
 \qquad \lvert\omega\rvert=1,\quad 0\leq a<1.
\]
The next lemma is independent of numerical ranges.  It is the step
that converts simultaneous estimates over the rank-one stretches
\eqref{eq:rank-one-stretch} into a sharper norm bound for the original
matrix.

\begin{lemma}[Rank-one stretch extraction]
\label{lem:orbit-extraction}
Let $n\geq2$, let $\kappa\geq1$ and $K\in\mathbb{R}$, and let $X\in\Mn$ satisfy
$\sigma(X)\subseteq\overline\D$.  Suppose that
\begin{equation}
 \lVert S_v\varphi_{\omega,a}(X)S_v^{-1}\rVert\leq K
 \label{eq:orbit-extraction-hypothesis}
\end{equation}
for every unit vector $v$, every $\lvert\omega\rvert=1$, and every
$0\leq a<1$.  Then
\begin{equation}
 \lVert X\rVert\leq\max\left\{1,\frac K\kappa\right\}.
 \label{eq:orbit-extraction-conclusion}
\end{equation}
\end{lemma}

\begin{proof}
Put $t=\lVert X\rVert$.  The result is immediate if $t\leq1$, so
assume $t>1$.  Choose unit vectors $x,y$ such that
\[
 Xx=ty.
\]
Let $c_0=\langle x,y\rangle$, put
$\omega=c_0/\lvert c_0\rvert$ when $c_0\neq0$, and put $\omega=1$
otherwise.  Set
\[
 \widetilde y=\omega y.
\]
Then $\lvert\omega\rvert=1$,
$\omega Xx=t\widetilde y$, and
$c:=\langle x,\widetilde y\rangle=\lvert c_0\rvert\in[0,1]$.
In fact $c<1$, because $c=1$ would give $\widetilde y=x$ and
$Xx=t\overline\omega x$, producing an eigenvalue of $X$ of modulus
$t>1$, whereas $\sigma(X)\subseteq\overline\D$.

If $c=0$, put $a=0$.  If $c>0$, consider
\[
 F(a)=tc(1+a^2)-a(1+t^2).
\]
Since
\[
 F(0)=tc>0,\qquad
 F(1)=2tc-(1+t^2)\leq-(t-1)^2<0,
\]
there is an $a\in(0,1)$ such that
\begin{equation}
 tc(1+a^2)=a(1+t^2).
 \label{eq:orthogonality-parameter}
\end{equation}
Define
\[
 u=(I-a\omega X)x=x-at\,\widetilde y,
 \qquad
 w=(\omega X-aI)x=t\,\widetilde y-ax.
\]
The matrix $I-a\omega X$ is invertible because
$\sigma(X)\subseteq\overline\D$ and $a<1$, so $u\neq0$ and
\[
 \varphi_{\omega,a}(X)u=w.
\]
Also $w\neq0$, since $w=0$ would imply
$t=\lVert Xx\rVert_2=a<1$.  Equation
\eqref{eq:orthogonality-parameter} gives
\[
 \langle u,w\rangle
 =tc(1+a^2)-a(1+t^2)=0.
\]
It also yields
\begin{align}
 \lVert w\rVert_2^2-t^2\lVert u\rVert_2^2
 &=a(t^2-1)\bigl(2tc-a(1+t^2)\bigr)\notag\\
 &=atc(t^2-1)(1-a^2)\geq0.
 \label{eq:automorphism-amplification}
\end{align}
Consequently,
\[
 \frac{\lVert w\rVert_2}{\lVert u\rVert_2}\geq t.
\]

Put $v=w/\lVert w\rVert_2$.  Since $u\perp w$, the stretch $S_v$
satisfies $S_v^{-1}u=u$ and $S_vw=\kappa w$.  Testing
$S_v\varphi_{\omega,a}(X)S_v^{-1}$ on $u$ gives
\[
 K\geq\lVert S_v\varphi_{\omega,a}(X)S_v^{-1}\rVert
 \geq\kappa\frac{\lVert w\rVert_2}{\lVert u\rVert_2}
 \geq\kappa t.
\]
Thus every case with $t>1$ has $t\leq K/\kappa$, proving
\eqref{eq:orbit-extraction-conclusion}.
\end{proof}

\section{Transfer to scaled \texorpdfstring{$q$}{q}-numerical ranges}
\label{sec:main-proof}

We now combine the geometry and extraction into a general transfer
principle.

\begin{theorem}[Spectral-constant transfer]
\label{thm:constant-transfer}
Let $K\geq1$ have the following property: for every square matrix $B$
and every rational function $h$ pole-free on $W(B)$,
\begin{equation}
 \lVert h(B)\rVert
 \leq K\max_{z\in W(B)}\lvert h(z)\rvert.
 \label{eq:base-numerical-range-bound}
\end{equation}
Let $n\geq2$, $A\in\Mn$, $0<\lvert q\rvert\leq1$, and let $f$ be
a rational function pole-free on $\Omega_q(A)$.  With
\[
 r=\lvert q\rvert,\qquad
 \kappa=\frac{1+\sqrt{1-r^2}}r,
\]
one has
\begin{equation}
 \lVert f(A)\rVert
 \leq
 \max\left\{1,\frac K\kappa\right\}
 \max_{z\in\Omega_q(A)}\lvert f(z)\rvert.
 \label{eq:transferred-bound}
\end{equation}
\end{theorem}

\begin{proof}
By \eqref{eq:phase-reduction}, it is enough to use $\Omega_r(A)$.
Put
\[
 M=\max_{z\in\Omega_r(A)}\lvert f(z)\rvert,\qquad
 M_\varepsilon=M+\varepsilon,\qquad
 X_\varepsilon=\frac{f(A)}{M_\varepsilon},
\]
where $\varepsilon>0$.  The maximum exists because $\Omega_r(A)$ is
compact and $f$ is continuous there.  Since
\[
 \sigma(A)\subseteq W(A)\subseteq\Omega_r(A),
\]
the rational spectral mapping theorem gives
$\sigma(X_\varepsilon)\subset\D$.

Fix a unit vector $v$ and put $B=S_vAS_v^{-1}$.
Lemma~\ref{lem:rank-one-stretch} gives
\[
 W(B)\subseteq\Omega_r(A).
\]
For $\lvert\omega\rvert=1$ and $0\leq a<1$, define the rational
function
\[
 h(z)=\varphi_{\omega,a}
      \left(\frac{f(z)}{M_\varepsilon}\right).
\]
The function $f$ is pole-free on $W(B)$, and the strict inequality
$M/M_\varepsilon<1$ keeps its normalized values in $\D$.  For every
$z\in W(B)$,
\[
 \left\lvert1-\frac{a\omega f(z)}{M_\varepsilon}\right\rvert
 \geq1-\frac{aM}{M_\varepsilon}>0.
\]
Thus $h$ is pole-free on $W(B)$ and $\lvert h\rvert\leq1$ there.
Applying
\eqref{eq:base-numerical-range-bound} and using similarity covariance
and composition covariance of the rational functional calculus, we
obtain
\[
 \lVert S_v\varphi_{\omega,a}(X_\varepsilon)S_v^{-1}\rVert
 =
 \left\|\varphi_{\omega,a}
  \left(\frac{f(B)}{M_\varepsilon}\right)\right\|
 \leq K.
\]
Lemma~\ref{lem:orbit-extraction} therefore gives
\[
 \frac{\lVert f(A)\rVert}{M_\varepsilon}
 \leq\max\left\{1,\frac K\kappa\right\}.
\]
Letting $\varepsilon\downarrow0$ proves
\eqref{eq:transferred-bound}.
\end{proof}

The sharp numerical-range theorem gives
\eqref{eq:base-numerical-range-bound} with $K=2$ for every rational
$h$ pole-free on $W(B)$ \cite{Jin2026}.

\begin{proof}[Proof of Theorem~\ref{thm:main}]
Apply Theorem~\ref{thm:constant-transfer} with $K=2$ and use
\eqref{eq:kappa-identities}.  This proves \eqref{eq:main-bound}.
\end{proof}

We next establish optimality of both branches of $C_q$.

\begin{proposition}[Sharpness]
\label{prop:sharpness}
For every fixed $q$ with $0<\lvert q\rvert\leq1$, put
$r=\lvert q\rvert$.  The constant
\[
 C_q=\max\left\{1,\frac2\kappa\right\},
 \qquad
 \kappa=\frac{1+\sqrt{1-r^2}}r,
\]
is the least constant $C$ such that
\[
 \lVert f(A)\rVert
 \leq C\max_{z\in\Omega_q(A)}\lvert f(z)\rvert
\]
for every $A\in\mathbb M_2(\C)$ and every rational function $f$
pole-free on $\Omega_q(A)$.
\end{proposition}

\begin{proof}
Theorem~\ref{thm:main} shows that $C_q$ is admissible.  By
\eqref{eq:phase-reduction}, $\Omega_q(A)=\Omega_r(A)$.  The constant
function $f\equiv1$ gives the lower bound $C\geq1$.

For the other branch, let
\[
 J=\begin{pmatrix}0&1\\0&0\end{pmatrix}.
\]
For a unit vector $x=(\xi,\eta)$, one has
\[
 \langle Jx,x\rangle=\overline\xi\eta,
 \qquad
 \rho_J(x)=\lvert\eta\rvert^2.
\]
Put $\mu=\lvert\xi\rvert$ and $\nu=\lvert\eta\rvert$, so that
$\mu^2+\nu^2=1$.  The identities in \eqref{eq:kappa-identities} give
\[
 \frac{\sqrt{1-r^2}}r=\frac{\kappa^2-1}{2\kappa}.
\]
Therefore every point $z$ in the disk indexed by $x$ in
\eqref{eq:q-disk-formula} satisfies
\begin{align*}
 \lvert z\rvert
 &\leq \mu\nu+\frac{\kappa^2-1}{2\kappa}\nu^2
 \leq\frac\kappa2.
\end{align*}
The second inequality follows from
\begin{equation}
 \frac\kappa2-
 \left(\mu\nu+\frac{\kappa^2-1}{2\kappa}\nu^2\right)
 =\frac{(\kappa\mu-\nu)^2}{2\kappa}\geq0.
 \label{eq:jordan-disk-bound}
\end{equation}
Equality is attained by taking the unit vector $x=(\mu,\nu)$ with
positive real coordinates
\[
 \mu=\frac1{\sqrt{\kappa^2+1}},
 \qquad
 \nu=\frac\kappa{\sqrt{\kappa^2+1}}.
\]
For this vector, the positive boundary point is $z=\kappa/2$.  Hence
\[
 \max_{z\in\Omega_r(J)}\lvert z\rvert=\frac\kappa2.
\]
Since $\lVert J\rVert=1$, taking the polynomial $f(z)=z$ yields
\[
 \frac{\lVert f(J)\rVert}
 {\max_{z\in\Omega_r(J)}\lvert f(z)\rvert}
 =\frac2\kappa.
\]
Thus every admissible $C$ satisfies $C\geq2/\kappa$.  Together with
$C\geq1$, this proves the claim.
\end{proof}

\section{Consequences and discussion}
\label{sec:discussion}

The transfer theorem accesses the base numerical-range estimate solely
through \eqref{eq:base-numerical-range-bound}, applied to the rank-one
stretches \eqref{eq:rank-one-stretch} selected by the extraction
mechanism.  The disk-automorphism family in
Lemma~\ref{lem:orbit-extraction} supplies the orthogonal input-output
pair used in \eqref{eq:automorphism-amplification}; the identity
function is one member of this family.

More generally, Theorem~\ref{thm:constant-transfer} turns every
universal numerical-range spectral constant into an explicit constant
for scaled $q$-numerical ranges.

For the sharp value $K=2$, the transfer constant equals $1$ precisely
when $\kappa\geq2$, or equivalently when
$0<\lvert q\rvert\leq4/5$.  At $\lvert q\rvert=1$, one has
$\kappa=1$ and the argument reduces to the sharp numerical-range
estimate.  The definition of the scaled range uses
$0<\lvert q\rvert\leq1$.  When $n\geq2$, unit vectors satisfying
$y^*x=q$ exist for every such $q$.

\section*{AI-assistance disclosure}

OpenAI ChatGPT contributed the rank-one stretch argument and the
disk-automorphism construction in Lemma~\ref{lem:orbit-extraction}.
It also assisted with exposition, bibliography, typesetting, and Lean
formalization.  The human author independently checked the mathematical
derivations and assumes full responsibility for the correctness of all
statements and for the integrity and accuracy of all citations.

\reftitle{References}
\bibliography{q_numerical_range_spectral_set}

\end{document}
